\documentclass[english,10pt,a4paper,draft,onecolumn]{article}
\usepackage[T1]{fontenc}
\usepackage[width=17cm]{geometry}
\usepackage{babel}
\usepackage{enumitem}
\title{Specification of the Theory-Native SPL Domain Specific Language}
\author{Ulrich Eisenecker, Yannis Hübenthal}
\begin{document}
\maketitle

\section{Purpose}
This document specifies the domain specific language accepted by the current theory-native SPL implementation. The DSL describes a Boolean base feature set, a conjunction of logical constraints, optional availability information about valid systems, and an optional output-side category filter.

\section{Token Types}
Token types are \emph{comment}, \emph{keyword}, \emph{name}, \emph{value}, and symbolic operators.
\begin{itemize}[leftmargin=*]
\item A comment begins with \emph{//} and must be followed by white space (space or tab). It extends to the next newline character.
\item Keywords are listed below.
\item A name is a sequence of non-white-space characters and is terminated by white space or punctuation.
\item A value is an integer.
\item Symbolic operators are \emph{!}, \emph{\&\&}, \emph{||}, \emph{->}, \emph{<->}, \emph{(}, \emph{)}, \emph{[}, \emph{]}, \emph{,}, and the comparison operators.
\end{itemize}

\section{Keywords}
Keywords are \emph{ASSERT}, \emph{AVAILABLE\_SYSTEMS}, \emph{CATEGORY\_FILTER}, \emph{CONSTRAINTS}, \emph{COUNT}, \emph{END}, \emph{FALSE}, \emph{FEATURES}, \emph{IN}, \emph{MISSING\_SYSTEMS}, \emph{SPL}, and \emph{TRUE}. Keywords are not case-sensitive.

\section{SPL}
\texttt{SPL} has cardinality 0..1. It specifies the optional name of the software product line. The name must not contain white space.

\section{FEATURES}
\texttt{FEATURES} has cardinality 1. It specifies the number of Boolean base features followed by exactly that many feature names.

Example:
\begin{quote}
\texttt{FEATURES 3 fa fb fc}
\end{quote}

\section{CONSTRAINTS}
\texttt{CONSTRAINTS} has cardinality 0..1. It contains zero or more \texttt{ASSERT} statements and is concluded by \texttt{END}. The semantics is the conjunction of all asserted formulas.

If \texttt{CONSTRAINTS} is omitted, there are no constraints. If \texttt{CONSTRAINTS} is present but contains no \texttt{ASSERT} statements, there are also no constraints.

Example:
\begin{quote}
\texttt{CONSTRAINTS\\
ASSERT fa -> fb\\
ASSERT !(fb \&\& fc)\\
END}
\end{quote}

\subsection{Supported formulas}
Formulas may contain:
\begin{itemize}[leftmargin=*]
\item feature variables
\item \texttt{TRUE} and \texttt{FALSE}
\item negation \texttt{!}
\item conjunction \texttt{\&\&}
\item disjunction \texttt{||}
\item implication \texttt{->}
\item equivalence \texttt{<->}
\item parentheses
\item cardinality expressions of the form \texttt{COUNT(...)}
\end{itemize}

The operator precedence is:
\begin{enumerate}[leftmargin=*]
\item \texttt{!}
\item \texttt{\&\&}
\item \texttt{||}
\item \texttt{->}
\item \texttt{<->}
\end{enumerate}

\subsection{Cardinality expressions}
A cardinality expression counts how many of the listed features are present in a system.

Supported forms are:
\begin{itemize}[leftmargin=*]
\item \texttt{COUNT(f1, f2, ..., fn) = k}
\item \texttt{COUNT(f1, f2, ..., fn) != k}
\item \texttt{COUNT(f1, f2, ..., fn) < k}
\item \texttt{COUNT(f1, f2, ..., fn) <= k}
\item \texttt{COUNT(f1, f2, ..., fn) > k}
\item \texttt{COUNT(f1, f2, ..., fn) >= k}
\item \texttt{COUNT(f1, f2, ..., fn) IN [a,b]}
\end{itemize}

\section{AVAILABLE\_SYSTEMS}
\texttt{AVAILABLE\_SYSTEMS} has cardinality 0..1. It specifies a set of valid systems known to be available. Every listed system must satisfy the conjunction of all asserted formulas.

If this block is present, it contributes to the explicit partition of the valid-system space.

\section{MISSING\_SYSTEMS}
\texttt{MISSING\_SYSTEMS} has cardinality 0..1. It specifies a set of valid systems known to be missing. Every listed system must satisfy the conjunction of all asserted formulas.

If this block is present, it contributes to the explicit partition of the valid-system space.

\section{Availability consistency}
A specification is invalid if:
\begin{itemize}[leftmargin=*]
\item a listed available system is not valid,
\item a listed missing system is not valid,
\item a system is listed as both available and missing, or
\item both blocks are present and their union is not exactly the valid-system set.
\end{itemize}

If both blocks are present, they must form a disjoint partition of the valid-system set. If only one of the two blocks is present, the other one is derived by complement with respect to the valid-system set. If neither block is present, all valid systems are treated as available.

\section{CATEGORY\_FILTER}
\texttt{CATEGORY\_FILTER} has cardinality 0..1. It is an optional output filter over the theory-native categories:
\begin{quote}
\texttt{Impossible Core Fundamental Not And Or Complex}
\end{quote}

\texttt{CATEGORY\_FILTER} affects reporting and post-classification selection only. It does not change valid systems, the generated feature registry, logical simplification, or category derivation.

\section{Canonical example files}
The repository keeps four paper-aligned example files as canonical references:
\begin{itemize}[leftmargin=*]
\item \texttt{spl\_specs/new\_dsl/valid/unconstrained\_2f.spl}
\item \texttt{spl\_specs/new\_dsl/valid/fx\_implies\_fy.spl}
\item \texttt{spl\_specs/new\_dsl/valid/fx\_implies\_fy\_missing\_s0.spl}
\item \texttt{spl\_specs/new\_dsl/valid/factorial\_paper.spl}
\end{itemize}
Documentation-facing copies of these examples are stored under
\texttt{external\_documentation/examples\_new\_dsl/} and should remain byte-for-byte aligned with the canonical fixtures.

\section{System semantics}
Let \(F\) be the Boolean base feature set declared by \texttt{FEATURES}. The candidate system space is \(P(F)\). The valid system space \(S_{valid}\) contains exactly those systems that satisfy the conjunction of all asserted formulas.

If no constraints are specified, then \(S_{valid} = P(F)\).

\section{Theory feature semantics}
The DSL does not enumerate theory features directly. Theory features are derived from the valid-system space, and their categories are derived from the simplified logical term. The DSL therefore contains no syntax for legacy generation controls.

\section{Removed legacy constructs}
The DSL contains no syntax for \texttt{MAX\_RANK}, legacy \texttt{CATEGORIES}, or \texttt{EXPLICIT\_DEPENDENCIES}. These constructs are outside the active theory-native language.

\end{document}
